
\documentclass{article} % For LaTeX2e
\usepackage{iclr2027_conference,times}

% Optional math commands from https://github.com/goodfeli/dlbook_notation.
\input{math_commands.tex}

\usepackage{hyperref}
\usepackage{url}
\usepackage{booktabs}
\usepackage{enumitem}


\title{What Actually Fixes an LLM Verifier --- And Why Nothing Else Does}

% The style file substitutes "Anonymous authors / Paper under double-blind review"
% for everything below until \iclrfinalcopy is uncommented.
\author{Anonymous authors \\
Paper under double-blind review}

%\iclrfinalcopy % Uncomment for camera-ready version, but NOT for submission.
\begin{document}


\maketitle

\begin{abstract}
LLM verifiers have a specific failure shape: they stay accurate at approving correct answers, but collapse at catching wrong ones---and five proposed fixes, tested across 64,800 verifications, all fail in that same shape.
Specifically, holding the exact error fixed, models approve their own wrong answers 11.9 percentage points more often than other models approve those same answers (across 6,969 strata spanning 779 distinct errors, 95\% CI [$+10.4, +13.5$])---and this gap holds regardless of how authorship is framed (\emph{you wrote this} / \emph{another model wrote this} / \emph{no framing}) or which verification strategy is used (direct verdict, chain-of-thought, or rubric).
All four verifiers---Qwen2.5-72B, DeepSeek-V3, Llama-3.3-70B, and Mistral-Nemo-12B---show the same asymmetric failure: they confirm correct answers 88.1\% of the time overall, but catch wrong ones only 46.1\% of the time, dropping to 26.4\% in the math domain---a pattern we term \emph{specificity collapse}.
To understand why, we ran five probes: style transfer (rewriting answers into another model's prose to test fingerprint recognition), self-recognition (asking models to identify their own answer from four candidates), structural self-preference (presenting two wrong answers side by side to test algorithmic favoritism), ensembling (testing whether 3-judge voting panels beat a single judge), and a frontier jury (two deliberating multi-model jury architectures without execution grounding).
The three mechanism probes all returned near-zero effects: style stripping moved false approval by at most $-1.5$\,pp ($p=0.13$); apparent self-recognition collapsed from $+22.2$\,pp to $+2.8$\,pp once we controlled for the model simply agreeing with its own prior conclusion ($p=0.73$); and structural favoritism, while real at 1.8--3.3$\times$, operates at absolute rates of only 3--13\%---far too small to account for an 11.9\,pp gap.
Ensembling and jury deliberation failed in the same direction as single judges, not independently of them, pointing to a shared underlying cause rather than a correctable surface bias.
What remains is \emph{belief persistence}: a model that reached a wrong conclusion during generation still believes that conclusion when reviewing the answer---the same mechanism as confirmation bias in human cognition, where we naturally seek evidence that confirms what we already believe rather than evidence that would disprove it.
The only intervention that breaks this pattern is an external ground-truth signal: when verifiers are shown execution test results, accuracy reaches 87.7\% on average---against 64.5\% for math and 57.2\% for science where no such signal exists---and a single grounded judge catches 93.3\% of errors compared to 76.7\% for the best deliberating multi-model jury without grounding.
No verifier-side fix---blinding, better prompting, ensemble aggregation, or multi-model deliberation---can substitute for a ground-truth signal; the reliability ceiling is set by what the environment can verify, not by the model reading it.
\end{abstract}

\section{Introduction}
\label{sec:intro}
% TODO: separate issue. Stub carries the label so cross-references resolve.

\section{Related Work}
\label{sec:related}
% TODO: separate issue. Stub carries the label so cross-references resolve.

\section{Experimental Setup}
\label{sec:setup}

\subsection{Tasks and Domains}
\label{sec:setup-tasks}

We draw 150 items from each of three benchmarks under a fixed sampling seed (42):
HumanEval+~\citep{liu2023evalplus} for code, MATH~\citep{hendrycks2021math} for
mathematics, and GPQA-Diamond~\citep{rein2024gpqa} for graduate-level
science.\footnote{HumanEval+/EvalPlus is distributed under Apache-2.0 (built on OpenAI's
HumanEval under MIT); MATH under the MIT License; GPQA-Diamond under CC BY-4.0, accessed
under its gated-release terms. All assets are used strictly for academic evaluation.}

The domains were chosen not for difficulty but for what the environment can verify.
Table~\ref{tab:domains} records the distinction the paper turns on. Correctness is
established by running the candidate in code, by matching a final answer in mathematics,
and by a multiple-choice letter in science, but only in code does any of that reach the
judge. Mathematics is therefore a checkable task presented to an unaided judge, separating \emph{grounding} (what the judge sees) from \emph{checkability} (what could be confirmed); if verifier reliability tracked checkability, math would behave like code.

For science, candidate options are shuffled per item under a per-item seed ($42 + \text{item index}$) and rendered as A through D before reaching models, so option position carries no answer signal. Items failing validity checks (non-four count, empty option, or duplicate text) are discarded before the 150-item sample is filled.

\begin{table}[t]
\caption{The three domains differ in whether a mechanical check exists and, separately, in
whether the judge is shown one. Only code verifications receive an environment signal.}
\label{tab:domains}
\begin{center}
\small
\begin{tabular}{llll}
\toprule
\bf Domain & \bf Benchmark & \bf Correctness established by & \bf Signal shown to judge \\
\midrule
Code    & HumanEval+   & Execution against the test harness     & Execution result block \\
Math    & MATH         & Final-answer extraction, numeric match & None \\
Science & GPQA-Diamond & Option letter; ambiguous parses fail    & None \\
\bottomrule
\end{tabular}
\end{center}
\end{table}

\subsection{Models}
\label{sec:setup-models}

The same four instruction-tuned models act as both generators and verifiers:
Qwen2.5-72B-Instruct, DeepSeek-V3, Llama-3.3-70B-Instruct-Turbo, and
Mistral-Nemo-Instruct-2407. All are served
through DeepInfra at temperature~0, with a 2500-token budget for generation and 2000 for
verification, with up to two retries on transport failure. The study is inference-only: the full grid consumes 51.8M prompt and 3.15M completion tokens across the 64{,}800 verification calls, alongside a CPU-only subprocess harness for code execution.

The pool deliberately spans a wide capability range, from a 12B dense model to a
671B-parameter mixture of experts. Section~\ref{sec:results-detect} exploits that spread to
ask whether stronger generators produce harder-to-catch errors;
Section~\ref{sec:limitations} notes it as a confound for any model-level ranking.

\subsection{The Verification Grid}
\label{sec:setup-grid}

Every model produces a candidate for every item, with no generation failures, giving 1{,}800
candidates. Every model then verifies every candidate. Each verification varies two factors
independently.

The \emph{authorship frame} tells the judge the answer is its own (\textsc{self}), another
model's (\textsc{other}), or says nothing (\textsc{neutral}). Prompts are identical apart
from that single sentence, and no model is ever named, so the frame carries a claim about
authorship and nothing else. The \emph{verification strategy} requests a bare verdict
(\textsc{direct}), a brief chain of thought before the verdict (\textsc{cot}), or a
rubric-scored assessment (\textsc{rubric}); the evaluation criteria themselves are held
constant across the three. Crossing four generators, four verifiers, three frames and three
strategies yields 144 evaluations per item, totaling 64{,}800 verification trials across 450 benchmark items and 1{,}800 unique candidates. To account for repeated evaluations over fixed items, all subsequent statistical tests and bootstrap confidence intervals are clustered at the item level.

The design's core property is that the told frame is crossed with true generator identity rather than confounded with it: comparing frames at fixed true authorship isolates the \emph{belief} of authorship, while comparing true authorship at a fixed frame isolates the \emph{fact} of it (\S\ref{sec:results-belief}).

One verification in four has a model grading its own candidate. Those 16{,}200 rows are the
object of study in Section~\ref{sec:results-belief}, and are excluded from every analysis
that compares judges to one another, namely detectability, oversight and ensembling, so
that self-preference cannot inflate a cross-judge comparison.

\subsection{Execution Grounding and Differential Fuzzing}
\label{sec:setup-grounding}

Code verifications, and only code verifications, carry a block reporting what happened when
the candidate was actually run against the benchmark's harness: exit status, exit code, the
number of assertions passed, captured output, and the final traceback frame on failure.
Runs exceeding a five-second budget are reported to the judge as a timeout rather than a
failure. The prompt states that the block is a runtime signal rather than the final word,
and explicitly permits the judge to mark passing code incorrect on logical grounds. That
permission is deliberate: it makes judge-initiated deviation from the environment signal
observable. Of the 21{,}600 code verifications, 1{,}458 (6.8\%) are such deviations, where
the judge rejected a candidate the harness had passed.

Adjudicating an override requires a check independent of both the harness and the judge, so
we fuzz each one differentially. A model outside the evaluated pool
(Gemma-2-27B-Instruct) proposes fifteen adversarial inputs aimed at boundary and edge
behavior; the candidate and the benchmark's reference solution are both run on them (median
15 input sets executed per override, maximum 83); and a second independent model
(WizardLM-2-8x22B) adjudicates any divergence, deciding which of the two outputs is correct.
We treat the oracle's verdict as a soft signal, since the oracle is itself unevaluated.
Each override lands in one of three resolved verdicts: \textbf{\texttt{NO\_DISCREPANCY}} (1{,}125 overrides, 77.2\%), where no divergence was found within the budget and the judge rejected code behaving like the reference; \textbf{\texttt{BUG\_CONFIRMED}} (107, 7.3\%), where a divergence was found and the oracle confirmed a genuine candidate bug; or \textbf{\texttt{REFERENCE\_BUG}} (11, 0.8\%), where the oracle judged the benchmark's own reference solution to be flawed.

The remaining 215 overrides (14.7\%) are unresolved: 168 (11.5\%) where the input-generation
step failed and 47 (3.2\%) where the fuzzer itself errored. These keep the harness verdict.

Two consequences carry into the rest of the paper. First, the fuzzer sits inside the
measurement loop rather than beside it. \texttt{BUG\_CONFIRMED} relabels a candidate as
incorrect, and \texttt{NO\_DISCREPANCY} and \texttt{REFERENCE\_BUG} relabel it as correct,
so every result is reported against both the raw harness labels and these
\emph{fuzz-adjusted} labels. Second, reference-solution bugs are rare but real: across every
override fuzzed, the oracle flagged the HumanEval+ reference itself on five items
(\texttt{HumanEval\_17}, \texttt{\_59}, \texttt{\_78}, \texttt{\_126} and
\texttt{\_134}), and the eleven \texttt{REFERENCE\_BUG} verdicts inside this grid fall on
four of those five. Even a gold-standard suite has blind spots, which is an argument for
grounding judges in execution rather than an argument against it.
Appendix~\ref{app:execution} gives the harness, its resource limits, and the fuzzer in full.

\subsection{Metrics}
\label{sec:setup-metrics}

Verifiers return JSON, and a verdict is read from a single boolean \texttt{is\_correct}
field. Of the 64{,}800 calls, 277 (0.43\%) did not return parseable JSON; for all 277 a
regex fallback recovered the verdict from the raw response
(Appendix~\ref{app:prompts}), so no call was discarded and every analysis below runs on the
full 64{,}800. The strict-parse failures are concentrated rather than diffuse: 262 come from
the smallest model (Mistral-Nemo) and 15 from Llama; 140 are \textsc{rubric} and 137
\textsc{cot}, with none from \textsc{direct}; 267 are math and 10 science, with none from
code.

We report the following item-level metrics throughout: (1)~\textbf{Accuracy} (\%): proportion of verdicts matching the ground-truth correctness label (correct $\to$ approved; incorrect $\to$ rejected); (2)~\textbf{Balanced accuracy}: the arithmetic mean of true-positive rate (TPR, sensitivity) and true-negative rate (TNR, specificity), $\tfrac{1}{2}(\mathrm{TPR}+\mathrm{TNR})$ (our \emph{primary metric}, as accuracy alone obscures specificity collapse); (3)~\textbf{False-positive rate (FPR)} / \textbf{false-approval rate}: the rate at which incorrect candidates are approved, $\mathrm{FP}/(\mathrm{FP}+\mathrm{TN})$; (4)~\textbf{False-negative rate (FNR)} / \textbf{false-rejection rate}: the rate at which correct candidates are rejected, $\mathrm{FN}/(\mathrm{FN}+\mathrm{TP})$; and (5)~\textbf{Item-adjusted approval rate} (Section~\ref{sec:results-detect} only): the within-item demeaned false-approval rate, removing item difficulty as a confound across generators.\footnote{Two senses of ``adjusted'' appear in this paper: \emph{fuzz-adjusted} refers to ground-truth relabeling (\S\ref{sec:setup-grounding}), while \emph{item-adjusted} refers to within-item demeaning (Appendix~\ref{app:capability}).}


\section{Results}
\label{sec:results}
% TODO: separate issue. Stub carries the label so cross-references resolve.

\subsection{Belief and Fact of Authorship}
\label{sec:results-belief}
% TODO: separate issue.

\subsection{Capability and Detectability}
\label{sec:results-detect}
% TODO: separate issue.

\section{Limitations and Future Directions}
\label{sec:limitations}
% TODO: separate issue.

\section{Conclusion}
\label{sec:conclusion}
% TODO: separate issue.

\subsection*{AI use statement}
% TODO: required by ICLR 2027, does not count toward the page limit.

\subsection*{Reproducibility statement}
% TODO: recommended by ICLR 2027, does not count toward the page limit.

\bibliography{refs}
\bibliographystyle{iclr2027_conference}

\appendix

\section{Execution Harness and Differential Fuzzer}
\label{app:execution}
% TODO: separate issue. Stub carries the label so cross-references resolve.

\section{Verification Prompts}
\label{app:prompts}
% TODO: separate issue.

\section{Capability and Grounding: Full Breakdown}
\label{app:capability}
% TODO: separate issue.

\end{document}
